\documentclass[11pt]{article}

\usepackage[utf8]{inputenc}
\usepackage[T1]{fontenc}
\usepackage{lmodern}
\usepackage{microtype}
\usepackage{amsmath, amssymb, amsthm}
\newtheorem{proposition}{Proposition}
\usepackage{geometry}
\geometry{a4paper, margin=1in}
\usepackage{xcolor}
\usepackage{hyperref}
\hypersetup{colorlinks=true, linkcolor=blue!50!black, citecolor=blue!50!black, urlcolor=blue!50!black}
\usepackage[backend=biber, style=numeric-comp, sorting=nyt]{biblatex}
\addbibresource{atree.bib}

\newcommand{\atree}{\texttt{atree}}

\title{The Atree Format: A Scalable Binary-Path Notation for\\ Ancestral Genealogies, with an Application to Lineage Matching}
\author{Vasili Gavrilov\\
\small Independent researcher\\
\small ORCID: \href{https://orcid.org/0009-0007-9371-5994}{0009-0007-9371-5994}}
\date{\today}

\begin{document}
\maketitle

\begin{abstract}
We describe \atree, a compact, human-readable, plain-text notation for \emph{ascending}
genealogies (a person's direct ancestors). Each ancestor is named by the path of ancestral sexes
from the subject, \texttt{M} for father, \texttt{F} for mother, so \texttt{MF} is the father's
mother. This path is exactly the classical Ahnentafel (Sosa--Stradonitz) number written in binary,
reading each $0$ as father and each $1$ as mother. A decimal Ahnentafel number doubles every
generation and soon grows unwieldy; an \atree{} path instead has one letter per generation, so its
length is simply how far back one is looking, and each line stands on its own, needing no surrounding
file, so any subset can be shared freely. We give the grammar, exact and reversible conversions to
and from Ahnentafel, tool support (GEDCOM import; an HIR/centimorgan calculator), and the limits
(ascending trees only; when the same ancestor recurs through cousin marriage, ``pedigree collapse'',
the tree becomes a graph). The main application is \emph{cross-person lineage matching}: when two
people's \atree{} paths share a tail, the names along it sound alike, and any Y/mtDNA haplogroups
agree, that is a measurable, spelling-robust signal that the two trees meet at the same ancestor.
About $37$ yes/no choices ($37$ bits) suffice to single out any human who has ever lived, so
identifying a person is itself a form of compression.
\end{abstract}

\noindent\textbf{Keywords:} genealogy; Ahnentafel; Sosa--Stradonitz; binary tree; ancestral path;
record linkage; phonetic name matching; haplogroup; data compression; Kolmogorov complexity.

\section{Introduction}
\label{sec:intro}
Personal genealogies must be stored and \emph{exchanged} as durable, human- and machine-readable
text. The classical Ahnentafel numbering is compact but its decimal identifiers grow exponentially
with depth and are opaque as paths. We present \atree: a sex-path identifier that is the natural
coordinate of the ancestor binary tree, fixing both the scalability and the self-containment
weaknesses of decimal Ahnentafel.\footnote{The notation underlying this paper arose earlier, in discussion on the molgen.org
genealogy forum and during the author's DNA/genealogy tool work of 2009--2011 (\texttt{aisconvert}
\textcopyright~2009 \cite{aisconvert}; \texttt{aisgedcom}/\texttt{AtreeParser} \textcopyright~2011
\cite{aisgedcom}); the binary-tree genealogy format itself is listed by the author as proposed in
2007. He named it ``\atree'' and first published the specification as a LiveJournal post on
7~November~2011 \cite{lj_atree}, adding a Wikipedia section the next day \cite{wiki_gns}. This paper
organizes and extends that material.}
Contributions: (i)~the \atree{} notation and self-contained plain-text format (building on the older
binary ancestor path of Anderson~\cite{anderson1995}); (ii)~a formal
grammar (with reference converters); (iii)~bijective conversions;
(iv)~a cross-person \emph{lineage-matching} method (combining structural, phonetic, and genetic evidence),
with an information-theoretic analysis; and (v)~a statement of limitations.

\section{Background: Ahnentafel / Sosa--Stradonitz}
\label{sec:background}
The Ahnentafel system \cite{eytzinger1590, wiki_ahnentafel} numbers a subject $1$, and assigns
person $n$'s father $2n$ and mother $2n{+}1$. It is credited to Eytzinger (1590) and takes its
common name ``Sosa--Stradonitz'' from Jer\'onimo de Sosa (1676) \cite{sosa1676} and Stephan
Kekul\'e von Stradonitz (1898) \cite{kekule1898}, who popularised it. This is the same numbering
computer science uses to keep a binary tree in a plain array (a node's two children at positions
$2n$ and $2n{+}1$), and listing ancestors in number order simply walks the family tree generation by
generation \cite{knuth_taocp3, clrs}. Writing $n$ in binary, dropping the leading $1$, then reading
each bit as father ($0$) or mother ($1$) gives the genealogical path; \atree{} renames those bits
\texttt{M}/\texttt{F}. This binary/letter-path
reading of Ahnentafel, writing an ancestor as the string of father/mother steps given by the binary
digits of its number, was already set out by Anderson~\cite{anderson1995} (in \texttt{f}/\texttt{m}
form); \atree{} adopts the \texttt{M}/\texttt{F} convention and builds a self-contained exchange
format upon it.

\section{The \atree{} notation}
\label{sec:notation}
An \atree{} identifier is the \texttt{M}/\texttt{F} path from the subject to an ancestor; the first
character encodes the subject's own sex. Formally, $\mathrm{atree}(n)=s\,\mu(b(n))$, where $s$ is
the subject's sex letter, $b(n)$ is the binary expansion of $n$ without its leading~$1$, and $\mu$
maps $0\mapsto\mathtt{M}$ and $1\mapsto\mathtt{F}$.
Each line carries its full path, so any subset of lines is valid and shareable on its own.
\textbf{No surrogate keys.} The path is a \emph{natural} key: it is derived from the data it names and
it states the position it denotes, so an identifier can be read rather than resolved. This is the point
of departure from both neighbours. GEDCOM assigns arbitrary surrogate identifiers
(\texttt{@I1@}) that carry no information and must be dereferenced. Ahnentafel numbers are derived
rather than arbitrary, but opaque: recovering the path from the number takes arithmetic, and the number
grows geometrically with depth. \atree{} introduces no surrogate identifier at any level.

A single \atree{} starts from one subject and lists only that subject's ancestors; a family or a
whole database is therefore a \emph{collection} of \atree{}s, one per person. Because paths are
relative to their subject, the same path string occurs in every tree and denotes a different person in
each, so a collection needs one identifier per \emph{tree} to say whose ancestors a block of lines
describes. That is the only identifier the format requires, and it scales with the number of subjects
rather than the number of people: $N$ trees need $N$ such labels, not one per ancestor. Joining two
collections is therefore a matter of keeping the subjects apart, not of reconciling per-person keys,
and the matching of Section~\ref{sec:matching} then works across trees on the paths themselves. Siblings are simply
different starting points whose ancestor paths coincide from the parents upward, so their lists share
the same tails, the overlap that the matching of Section~\ref{sec:matching} uses.

\section{Grammar}
\label{sec:grammar}
The grammar is BNF with POSIX character classes and regular-expression repetition. One rule is
semantic rather than syntactic: a parenthesized group \texttt{(\dots)} is a \emph{date} when its
contents parse as a date and it is \emph{not} enclosed by the surname slashes \texttt{/\dots/} (in
which case it is a maiden name).

\begin{verbatim}
<atree>       ::= ( <line> <eol> )*
<line>        ::= <blank> | <comment_line> | <person>
<blank>       ::= <ws>*
<comment_line> ::= <ws>* <comment>
<comment>     ::= "#" [:print:]*

<person>      ::= <path> <haps> <namefield> <datefield>
                  <locfield> <commentfield>
<haps>        ::= ( <ws>+ <haplogroup> )*
<namefield>   ::= ( <ws>+ ","? <ws>* <names> )?
<datefield>   ::= ( <ws>* <dates> )?
<locfield>    ::= ( <sep> <location> )?
<commentfield> ::= ( <ws>* <comment> )?
<sep>         ::= <ws>+ ","? <ws>* | "," <ws>*

<path>        ::= [MF]+     ; 1st letter = subject's sex;
                            ; rest = father M / mother F
<haplogroup>  ::= "[" ( <hapsys> ":" )? <hapname> "]"
<hapsys>      ::= "Y" | "mt"   ; the two trees share clade letters
<hapname>     ::= [:alnum:] ( [:alnum:] | "-" | "." )*
                            ; stable SNP form, e.g. R-M269, N-M178
<names>       ::= <name_text>? ( <ws>* <surname>
                  ( <ws>* <name_text> )? )?
<name_text>   ::= <word> ( <ws>+ <word> )*
<word>        ::= <name_char>+
<name_char>   ::= [:alnum:] | "'" | "-" | "." | "?"
<surname>     ::= "/" <name_text>? <maiden>? "/"
<maiden>      ::= "(" <name_text>? ")"

<dates>       ::= "(" <ws>* <date_or_range> <ws>* ")"
<date_or_range> ::= <gdate> ( <ws>* "-" <ws>* <gdate> )?
<gdate>       ::= <approx>? ( <full_date> | <short_date>
                  | <year> ) | "?"
<approx>      ::= ( "ABT" | "BEF" | "AFT" | "EST" ) <ws>+ | "~"
<full_date>   ::= <day> <ws>+ <month> <ws>+ <year>
<short_date>  ::= <month> <ws>+ <year>
<day>         ::= [:digit:]{1,2}
<year>        ::= [:digit:]{1,4}
<month>       ::= "JAN"|"FEB"|"MAR"|"APR"|"MAY"|"JUN"
                | "JUL"|"AUG"|"SEP"|"OCT"|"NOV"|"DEC"

<location>    ::= <loc_char>+
<loc_char>    ::= [:print:] except "#"
<ws>          ::= [:blank:]
<eol>         ::= "\n"
\end{verbatim}

\section{Line syntax and examples}
\label{sec:examples}
Names precede an optional surname in \texttt{/\dots/} (GEDCOM style) with a maiden name in
parentheses; \texttt{//} and \texttt{?} mark unknowns; dates use \texttt{ABT}/\texttt{BEF}/%
\texttt{AFT}/\texttt{EST} and ranges. A worked ascending tree (excerpt); each person is one line:

{\footnotesize\begin{verbatim}
M    [Y:N1c1] ...
MM   [Y:N1c1] ... /Gavrilov/
MF   [mt:H10a1] Zoya Vasilyevna /Gavrilova (Guseva)/ (20 SEP 1941 - 20 OCT 1984)
MFM  Vasiliy Nikolayevich /Gusev/ (~1914 - 1963)
MFF  [mt:H10a1] Anastasiya Ivanovna /Guseva (Marinova)/ (20 DEC 1913 - 03 JUN 2006)
MFMF Evdokiya Nikolayevna /Guseva (?)/ (?), Milna, daughter of merchant
\end{verbatim}}

\section{Conversions}
\label{sec:conversions}
\begin{proposition}\label{prop:bij}
For a fixed subject sex, $\mathrm{atree}(\cdot)$ is a bijection between Ahnentafel numbers $n\ge 1$
and \atree{} paths.
\end{proposition}
\begin{proof}
Writing $n$ in binary and dropping the leading $1$ is a bijection from $\{n\ge 1\}$ onto
$\{0,1\}^{*}$; the relabelling $0\mapsto\mathtt{M}$, $1\mapsto\mathtt{F}$ is a bijection onto
$\{\mathtt{M},\mathtt{F}\}^{*}$; and the subject-sex letter is prepended (and, in reverse,
normalised) invertibly. The two converters below realise both directions.
\end{proof}
The reference converters are two short shell scripts,
\texttt{atree2s.sh} (Ahnentafel\,$\to$\,\atree) and \texttt{s2atree.sh} (\atree\,$\to$\,Ahnentafel),
reproduced in Appendix~\ref{app:scripts}; they are preferred over the Java toolkit for portability,
e.g.\ \texttt{s2atree.sh MFFMMMMFFFMMF} yields \texttt{7225}.

\begin{center}
\begin{tabular}{lcccc}
\hline
Person & Ahnentafel & Binary & \atree{} (subj.\ F) & \atree{} (subj.\ M)\\
\hline
Self          & 1 & 1   & F   & M\\
Father        & 2 & 10  & FM  & MM\\
Mother        & 3 & 11  & FF  & MF\\
Paternal GF   & 4 & 100 & FMM & MMM\\
Paternal GM   & 5 & 101 & FMF & MMF\\
Maternal GF   & 6 & 110 & FFM & MFM\\
Maternal GM   & 7 & 111 & FFF & MFF\\
\hline
\end{tabular}
\end{center}

\section{Properties}
\label{sec:properties}
\textbf{Record independence.} One line is one person, and it carries its full path, so it means the
same thing alone as in the file it came from. This is the property the rest follow from. Because no
line refers to another, a line can be questioned, corrected, re-sorted, dropped or sent on its own,
which matches how genealogical research actually proceeds: one ancestor at a time, provisionally, and
revised years later. The grammar supports the provisional state directly, with \texttt{?} and
\texttt{//} for unknown fields and \texttt{(?)} for an unknown maiden name (Appendix~\ref{app:template}),
so partial knowledge is written down rather than left absent. \textbf{Hand-repairability.} Having no
cross-references to keep consistent, the file can be edited in any text editor without a parser, and a
damaged file degrades line by line rather than as a whole. \textbf{Tool-independence.} A format is
durable to the extent that it can be read and repaired without its software, which is the practical
content of the digital-preservation argument for plain text. This work supplies its own evidence: the
2009--2011 Java tooling described in Section~\ref{sec:tooling} no longer builds unmodified, while the
files it produced remain readable and editable, and the reference converters of
Appendix~\ref{app:scripts} still run. \textbf{Short identifiers.} A path has one letter per
generation, so its length is just the generation distance (about $\log_2$ of the Ahnentafel number),
and there is no large number to overflow, unlike a decimal Ahnentafel. \textbf{Natural tree order.}
The letters give an ancestor's generation and position directly, so paths sort into family-tree order,
and being UTF-8 and one record per line the format is easy to compare (diff) and merge
(version-control friendly).

\section{Tooling and interoperability}
\label{sec:tooling}
\texttt{aisgedcom} \cite{aisgedcom} is the \atree{} reference implementation: a parser
(\texttt{AtreeParser}), GEDCOM$\leftrightarrow$\atree{} conversion, the canonical grammar and example
records, and a two-relatives relationship calculator. \texttt{aisconvert} \cite{aisconvert} (personal
genomics toolkit) provides the HIR/Half-IBD calculator, which measures shared
autosomal segments and reports centimorgans across autosomes 1--22. It computes sharing; it does not
infer relationships. The X chromosome is reported separately, since its sex-dependent transmission
traverses only a restricted subset of pedigree paths, so X centimorgans must not be pooled into a
total from which relationship distance is later inferred. The \atree{} term appears in third-party genealogy
discussion, for example a Gramps forum thread proposing a ``Dual Atree chart for cousins''
\cite{gramps_atree}. The wording there follows the Wikipedia section the author added in 2011
\cite{wiki_gns}, so the term's propagation is traceable to that edit and is not independent adoption.

\section{Limitations and open problems}
\label{sec:limitations}
\atree{} covers the \emph{ascending} (two-parent) tree; descendant trees branch many ways and are out
of scope. \textbf{Pedigree collapse} (the same ancestor appearing more than once through cousin
marriage) gives more than one path to a single real person: a path is unique as a \emph{position},
not as a \emph{person}, so the full pedigree is no longer a tree but a graph (one ancestor reached by
several paths), and storing it compactly becomes a matter of sharing the repeated sub-trees (a
compression problem). A full genealogy, covering all branches or a whole database, is therefore kept
as a \emph{collection} of \atree{}s, one ascending tree per subject (Section~\ref{sec:notation}); this
covers the whole graph at the cost of repeating shared ancestors across trees, the very repetition
that a collapse-aware, sub-tree-sharing encoding would remove.

\section{Application: cross-person lineage matching}
\label{sec:matching}
The central application is matching people who share ancestors, a long-standing problem in
genealogy: deciding when two records mean the same person. When persons $A$ and $B$ share a common
ancestor $X$, every ancestor \emph{above} $X$ is the same in both trees; so although their paths
\emph{to} $X$ differ (they start from different people), the \texttt{M}/\texttt{F} steps from $X$
upward, the tail of each path, are identical, and the names along that tail agree apart from
spelling. Four signals combine; the first two need only the written record: (i)~a shared \textbf{tail} of the path (the longer the
shared tail, the more generations in common, and this uses no names at all); and (ii)~\textbf{name
sound-alike} agreement \cite{russell1918, philips1990} to absorb spelling differences.
Phonetic coding is \emph{optional} and is performed by the matching software at comparison time, not
required by the format: names are stored as written, in whatever spelling is convenient to the user,
and the matcher applies Soundex/Metaphone (or any chosen scheme) when scoring. All normalisation is
deferred to the software, keeping the stored file human-convenient.
A missing or extra generation shifts the tail by a position or two, so the tails are compared by
approximate alignment, allowing an occasional inserted or dropped step, rather than by exact match.
Writing each path deep-ancestor-first and grouping paths by their shared opening avoids comparing
every pair: records with the same deep ancestry fall into the same group, so a match becomes a
lookup. In the language of record linkage, the shared tail is a new grouping key that does not depend
on names. The two extreme paths are exactly the DNA-trackable lines: all-\texttt{M}
($\mathtt{MMM}\dots$) is the patrilineal (Y-DNA) line, all-\texttt{F} ($\mathtt{FFF}\dots$) the
matrilineal (mtDNA) line, and the haplogroup field is the third signal, confirming or refuting a match along those two
edges alone. A fourth signal is independent of the written record altogether. Autosomal half-identical regions
(HIR), measured in centimorgans, are shared in an amount that falls with relationship distance, and the
\texttt{aisconvert} calculator \cite{aisconvert} computes them across autosomes 1--22 and X. Unlike
the other three, this one does not read the genealogy at all: it measures the genome, so it
\emph{tests} the relationship a shared tail asserts rather than adding to the assertion. A tail
implying a common ancestor $g$ generations back predicts an expected range of shared centimorgans;
agreement corroborates the paper trail, and disagreement localises an error in it. The check has a
range limit, since detectable autosomal sharing falls away beyond roughly six to eight generations, so
it validates only the recent end of a path. That is the same window in which, by the argument of
Section~\ref{sec:why}, the distinguishing information sits, so the two lines of reasoning bound the
useful depth from opposite directions.

Combining
the four pieces of evidence, the length of the shared tail, the name sound-alike, haplogroup
agreement on the uniparental edges, and autosomal centimorgan sharing, in a standard record-linkage
score \cite{fellegi1969} gives a match \emph{probability}. We present this as a
\emph{proposed} method: the false-positive model of Section~\ref{sec:why} motivates it analytically,
but empirical evaluation on a labeled genealogical corpus (precision and recall against known
relationships) remains future work.

\subsection{Why it works: far fewer real ancestors than possible paths}
\label{sec:why}
About $1.17\times10^{11}$ people have ever lived \cite{prb_everlived}; since
$\log_2(1.17\times10^{11})\approx 37$, roughly $37$ yes/no facts ($37$ bits) are enough to give
\emph{every} human who has ever lived a distinct label, close to the smallest code possible. (For
scale, $37$ generations span about a millennium; reaching an ancestor more distant than that needs a
longer path, but pedigree collapse, below, holds the number of \emph{distinct} ancestors far below
the naive count of $2^{L}$ possible paths.) An \atree{} path is also a label relative to its
starting person: the same ancestor has different paths from different descendants (and several when
collapse occurs), so the $37$-bit figure says only that a short label \emph{exists}; it does not
claim that an \atree{} path is by itself a single global ID for that person, and building such an ID
is the collapse-aware merge of Section~\ref{sec:conclusion}. There are $2^{L}$ possible paths of
length $L$, far more than the $\le 10^{11}$ people who could fill them (at $L{=}100$ that is
$2^{100}\approx 1.3\times10^{30}$ slots), so the paths actually used are very thinly spread. Two
unrelated lines share a tail of $k$ steps by chance with probability about $2^{-k}$ (taking the sexes
along a path as roughly even and independent; about $10^{-9}$ at $k{=}30$), and matching names make a
coincidence even less likely; a shared tail of $20$--$30$ letters with matching names is almost
certainly a real common ancestor. The paths in use are thinly spread \emph{because they are shared}:
pedigree collapse, together with the recent common ancestry of all living people
\cite{chang1999, rohde2004}, means deep ancestors are common to nearly everyone (the tails converge)
while people differ mainly in recent generations. So the information that tells people apart sits in
the recent end of the path (the first $\sim\!20$--$35$ letters), and the number of distinct
lines is about the population, $\sim\!\log_2(\text{population})\approx 37$ bits, not the naive $100$.
The huge count of possible paths is an artifact of writing the family graph out as a tree; the real
object is a small, shared graph, and identifying a person in it is again a form of compression.

\subsection{A scoring model}
\label{sec:scoring}
Record linkage scores a candidate pair by summing the log likelihood ratio of each piece of evidence
\cite{fellegi1969}; in bits, the terms add. Writing $R$ for ``the pair is related as the shared tail
asserts'':

\begin{center}
\renewcommand{\arraystretch}{1.2}
\begin{tabular}{@{}lll@{}}
\hline
Evidence & $\log_2 \mathrm{LR}$ (bits) & Where it is valid\\
\hline
Shared tail of $k$ steps      & $k$                       & recent range only (Sec.~\ref{sec:why})\\
Name agreement                & $\log_2(m/u)$             & $m,u$ estimated from a labelled corpus\\
Haplogroup concordant, freq.\ $f$ & $\log_2(1/f)$: $1.0$ at $f{=}0.5$, & all-\texttt{M} / all-\texttt{F} edges only\\
                              & $2.3$ at $f{=}0.2$, $4.3$ at $f{=}0.05$ & \\
Haplogroup discordant         & large negative (exclusion) & all-\texttt{M} / all-\texttt{F} edges only\\
Autosomal cM sharing          & likelihood over $g$        & $g \lesssim 6$--$8$ generations\\
\hline
\end{tabular}
\end{center}

\noindent\textbf{The tail term} follows from Section~\ref{sec:why}: two unrelated lines share a
$k$-step tail with probability about $2^{-k}$, while related ones share it by construction, so
concordance is worth about $k$ bits. This holds only in the range where the independence assumption
does: beyond roughly $20$--$30$ generations, universal recent common ancestry \cite{chang1999,
rohde2004} makes deep tails common to nearly everyone, and the term must not be read as $2^{k}$ there.
The model is therefore strongest exactly where Section~\ref{sec:why} locates the distinguishing
information, in the first $20$--$35$ letters.

\noindent\textbf{The haplogroup term is asymmetric,} and this is its main practical point.
Haplogroups are deep clades, unchanged over genealogical time, so concordance is expected of a true
relative but also of everyone else in the clade: its evidential value is only $\log_2(1/f)$, which for
a common haplogroup is one or two bits, an order of magnitude less than a ten-step tail. Discordance,
by contrast, is close to decisive, since a mismatch along an unbroken patriline or matriline refutes
the asserted path. The field is thus better read as an exclusion test than as confirmation, and it
speaks only about the two extreme paths; a mixed \texttt{M}/\texttt{F} tail receives no uniparental
evidence at all.

\noindent\textbf{The autosomal term is the only one that measures anything outside the record.}
Expected sharing falls with relationship distance, so an observed centimorgan total induces a
likelihood over $g$ that can be compared with the $g$ the tail asserts. Haplogroups cannot supply this,
having no genealogical-timescale clock; short tandem repeat haplotypes, which mutate at a rate per
marker per generation, could supply a second such estimate along the patriline.

We give the model rather than fitted parameters: $m$ and $u$, the haplogroup frequencies, and the
expected-sharing curve are all quantities the literature or a labelled corpus supplies, and estimating
them here without such a corpus would be a fit to nothing. The contribution is the decomposition and
its stated validity limits.

\section{Related work}
\label{sec:related}
The dominant genealogical interchange format is GEDCOM \cite{gedcom55} (and the later GEDCOM~X
\cite{gedcomx}):
record- and pointer-oriented, and general enough to represent arbitrary family graphs. \atree{} is
complementary rather than competing, a compact, line-oriented encoding specialized to the
\emph{ascending} tree, in which each line is self-sufficient. The difference is structural rather
than stylistic: GEDCOM identifies individuals by cross-reference, so a record's meaning depends on the
records it points to, and editing by hand means preserving referential integrity by hand. That design
is what makes GEDCOM general, and also what makes a single record unusable in isolation. In \atree{}
the path is the position, so there is nothing to dereference. When the format was designed
(2007--2011), genealogies were kept either in application-specific stores or in GEDCOM, so correcting
or exchanging part of a tree meant working through an application; independence from any application
was the motivating requirement. The author's own open-source tools
realize the surrounding interoperability: \texttt{aisgedcom} converts between GEDCOM and \atree{}
\cite{aisgedcom}; \texttt{aisconvert} \cite{aisconvert} parses 23andMe/FTDNA genomes and computes
HIR / centimorgan overlaps (with \texttt{ped2raw} \cite{ped2raw} for PED$\to$RAW conversion); and
\texttt{nid} \cite{nid} produces anonymized identifiers by cryptographic hashing, relevant to the
entity-resolution setting. Among genealogical numbering systems,
ascendant numbering is dominated by Ahnentafel / Sosa--Stradonitz \cite{eytzinger1590,
wiki_ahnentafel}, while descendant lines use the Henry, d'Aboville, de~Villiers--Pama, and Register
systems \cite{wiki_gns}; \atree{} is the binary (path) form of Ahnentafel and inherits its two-parent structure. The
binary/letter-path reading of Ahnentafel was described earlier by Anderson~\cite{anderson1995}, who
also observes that the charts of two people can be compared to determine their relationship. The
structural comparison is therefore Anderson's; the contribution of \atree{} over that precursor is the
named, self-contained exchange format with its grammar and reference implementations, and, for
matching, the scored combination of suffix, phonetic, uniparental and autosomal evidence with stated
validity limits. The
matching method of Section~\ref{sec:matching} builds on classical record linkage \cite{fellegi1969}
and phonetic name coding \cite{russell1918, philips1990}, contributing the \atree{} path suffix as a
compact, name-independent blocking key, complementary to autosomal-segment (HIR / centimorgan) and
uniparental (Y/mtDNA) evidence.

\section{Prior publication}
\label{sec:provenance}
The \atree{} specification, with its BNF grammar, was first published by the author on
7~November~2011 \cite{lj_atree}; the corresponding section of the Wikipedia article
\emph{Genealogical numbering systems} \cite{wiki_gns} was added the following day. The reference
implementations predate both (\texttt{aisconvert}, 2009; \texttt{aisgedcom}, 2011). The Wikipedia
section was written by the author, so it is disclosed here as provenance and is not cited as
independent corroboration. The present paper consolidates and extends that material: the grammar, the
bijection, the tooling description, and the lineage-matching analysis of
Section~\ref{sec:matching} appear here for the first time in a single reviewed account. Both 2011
sources are archived by the Internet Archive.\footnote{%
\url{http://web.archive.org/web/20260607185709/https://siberean.livejournal.com/14874.html} and
\url{http://web.archive.org/web/20260420170149/https://en.wikipedia.org/wiki/Genealogical_numbering_systems}.}

\section{Conclusion}
\label{sec:conclusion}
\atree{} encodes a direct ancestor as the \texttt{M}/\texttt{F} path from the subject, the binary
form of the Ahnentafel number, giving identifiers that are logarithmic in tree size,
self-sufficient, and plain-text (hence diffable, mergeable, and durable). It comes with a grammar,
bijective converters, and GEDCOM interoperation. Its main payoff is matching: a shared path suffix,
combined with phonetic name agreement and uniparental haplogroup concordance, gives a quantifiable,
name-robust signal for genealogical entity resolution, and $\sim\!37$ bits suffice to address any
human who has ever lived, identification as compression, a perspective developed in the author's
companion work on intelligence as the discovery of compressors~\cite{gavrilov_compressors}. Open
directions: a collapse-aware merge
that folds the tree-unfolding back to the underlying DAG, and joint structural/phonetic/genetic
scoring at scale.

\section*{Acknowledgements and attribution}
The bare sex-path notation, a string of \texttt{M}/\texttt{F} letters (e.g.\ \texttt{MMFFFMFM}), was
already in circulation in genealogy forum discussion, used by example and without a name; the author
met it there and thanks the correspondent who introduced it. The notation is pre-existing. The
earliest published description the author has located is Anderson's, in print by 1995
\cite{anderson1995}, which gives the same father/mother reading of the binary Ahnentafel expansion in
\texttt{f}/\texttt{m} form. The author was unaware of it: Anderson's two-page note appeared in a
regional genealogical society journal and did not surface in the general sources consulted at the
time, among them the Wikipedia survey of genealogical numbering systems, which then listed no binary
or path-based notation at all. The format described here was therefore developed without reference to
that precursor, which is a statement about the author's knowledge and not a claim of priority over
it. The author's distinct contributions are therefore \emph{not} the binary path
itself, but: the \textbf{name ``\atree''}; the \textbf{self-contained plain-text format and grammar};
the \textbf{reference implementations} as sole programmer; and the \textbf{lineage-matching analysis}
of Section~\ref{sec:matching}. The raw notation has a shared, pre-existing origin; the name,
formalization, implementation, and application are the author's.

The author's grounding in graph theory owes much to the advanced graph-theory course of the late
Professor Shimon Even (1935--2004) \cite{even_graph, even_comb}, whose student he was.

\section*{Code and data availability}
The \atree{} reference implementation and related tooling are open source: \texttt{aisgedcom}
(GEDCOM$\leftrightarrow$\atree{} conversion and \texttt{AtreeParser}) \cite{aisgedcom};
\texttt{aisconvert} (HIR / centimorgan calculator) \cite{aisconvert}; \texttt{ped2raw}
\cite{ped2raw}; and \texttt{nid} \cite{nid}. The two reference shell converters are reproduced in
full in Appendix~\ref{app:scripts}. No new datasets were used; the record in
Section~\ref{sec:examples} is an illustrative excerpt.
These tools were written in 2009--2011 and are no longer maintained; they are provided as historical
artifacts, for provenance and reference, and may not build or run unmodified on current systems. The
self-contained shell converters of Appendix~\ref{app:scripts}, by contrast, are current: run under a
POSIX shell, \texttt{s2atree.sh MFFMMMMFFFMMF} returns \texttt{7225} and \texttt{atree2s.sh 7225}
returns the same path, differing only in the leading subject-sex letter, which Ahnentafel does not
encode and which the script takes from its \texttt{my\_gender} setting (Proposition~\ref{prop:bij}).
Conversion was checked at every depth from~$1$ to~$200$: the path length matches the generation
distance, the alphabet is closed, and the round trip through \texttt{s2atree.sh} returns the original
number exactly.

\appendix

\section{A template / example file}
\label{app:template}
Person-field forms (each real line is prefixed by an \atree{} \texttt{M}/\texttt{F} path; the field
forms below are the part after the path). Surname in \texttt{/\dots/}, maiden name in parentheses,
\texttt{?} and \texttt{//} for unknowns, a date acts as the name/place delimiter, and \texttt{\#}
starts a comment.

\begin{verbatim}
FirstName
FirstName (1600)
FirstName (ABT 1600)
FirstName (14 AUG 1855-5 APR 1877)
FirstName /Surname/
FirstName /Surname/ (?-ABT 1800)
FirstName /Surname/ (AFT 1800-?)
/Surname/ (28 OCT 1852-BEF 1880)
FirstName /Surname (MaidenSurname)/ (1764-6 OCT 1813)
FirstName Another Name /Surname (MaidenSurname)/ (14 AUG 1855-5 APR 1877), Place of origin
FirstName // (1600)        # nothing but a date is known
//                         # nothing is known
/?/                        # surname refinement is the priority
FirstName /Surname (?)/    # maiden surname unknown
# a comma before the given name, as some languages / GEDCOM require:
/Surname/, AnotherName (1800), Place of origin
# if a place could be confused with a comma in a name, the date is the delimiter:
Names (?) Place of origin
\end{verbatim}

\section{Reference converter scripts}
\label{app:scripts}
The two reference converters (subject sex set at the top; \texttt{bc} does the base conversion).

\begin{verbatim}
#!/bin/sh
# atree2s.sh : convert an Ahnentafel number to an atree path  (POSIX sh).
# Usage: atree2s.sh <ahnentafel-number>

my_gender=F          # subject's sex: M or F
male=M
female=F

case "$1" in
    '')        echo "usage: $0 <ahnentafel-number>" >&2; exit 1 ;;
    *[!0-9]*)  echo "$0: '$1' is not a number"      >&2; exit 1 ;;
esac
case "$1" in
    *[!0]*)    : ;;
    *)         echo "$0: Ahnentafel numbers start at 1" >&2; exit 1 ;;
esac

# Ahnentafel number -> binary -> letters (0 = father M, 1 = mother F).
# bc folds long output at 70 columns with a backslash-newline;
# strip the continuations before touching the digits, or paths
# deeper than 67 generations are silently corrupted below.
binary=$(echo "obase=2; $1" | bc | tr -d '\\\n')
letters=$(printf '%s' "$binary" | tr '0' "$male" | tr '1' "$female")

# The leading bit denotes the subject; rewrite it as the subject's sex.
letters=$(printf '%s' "$letters" | sed "s/^./$my_gender/")

echo "$letters"
\end{verbatim}

\begin{verbatim}
#!/bin/sh
# s2atree.sh : convert an atree path to its Ahnentafel number  (POSIX sh).
# Usage: s2atree.sh <atree-path>

male=M
female=F

case "$1" in
    '')        echo "usage: $0 <atree-path>"        >&2; exit 1 ;;
    *[!MF]*)   echo "$0: '$1' is not an atree path" >&2; exit 1 ;;
esac

# The first letter is the subject; the Ahnentafel leading bit is always 1,
# so normalise it to F before mapping letters back to bits (M = 0, F = 1).
path="$female$(printf '%s' "$1" | cut -c2-)"
binary=$(printf '%s' "$path" | tr "$male" '0' | tr "$female" '1')

echo "ibase=2; $binary" | bc
\end{verbatim}

\printbibliography

\end{document}
